\documentclass[aspectratio=169,11pt]{beamer}

% ─── Clean minimal theme ────────────────────────────────────────────────────
\usetheme{default}
\usecolortheme{dove}
\usefonttheme{professionalfonts}
\setbeamertemplate{navigation symbols}{}
\setbeamertemplate{footline}[frame number]
\setbeamertemplate{frametitle}[default][left]

% Colors
\definecolor{accent}{HTML}{2171B5}
\definecolor{darktext}{HTML}{333333}
\definecolor{lightbg}{HTML}{F0F4F8}
\definecolor{greenaccent}{HTML}{2E7D32}
\definecolor{redaccent}{HTML}{C62828}
\definecolor{orangeaccent}{HTML}{E65100}
\setbeamercolor{title}{fg=accent}
\setbeamercolor{frametitle}{fg=accent}
\setbeamercolor{structure}{fg=accent}
\setbeamercolor{normal text}{fg=darktext}
\setbeamercolor{itemize item}{fg=accent}
\setbeamercolor{footline}{fg=gray}
\setbeamercolor{section in toc}{fg=accent}
\setbeamerfont{section in toc}{series=\bfseries}
\setbeamerfont{title}{series=\bfseries,size=\Large}
\setbeamerfont{frametitle}{series=\bfseries,size=\large}

% Packages
\usepackage{graphicx}
\usepackage{booktabs}
\usepackage{amsmath,amssymb}
\usepackage{mathtools}
\usepackage{tikz}
\usetikzlibrary{arrows.meta,positioning,calc,decorations.pathreplacing,shapes.geometric}
\usepackage{algorithm}
\usepackage{algorithmic}
\renewcommand{\algorithmiccomment}[1]{\hfill{\scriptsize\textcolor{gray}{$\triangleright$ #1}}}

% Invisible tcolorbox that scales content to fit frame body, centered.
\usepackage[fitting,skins]{tcolorbox}
\newtcolorbox{fitcenter}{%
  blanker,
  fit,
  fit basedim=11pt,
  halign=center,
  width=\textwidth,
  height=210pt,
}

\graphicspath{{../../figures/}{../../../figures/}}

\newcommand{\E}{\mathbb{E}}
\newcommand{\KL}{\mathrm{KL}}
\newcommand{\Var}{\mathrm{Var}}
\DeclareMathOperator{\Cov}{Cov}

\title{Mirror Descent Diffusion Policy}
\subtitle{Adaptive Inference-Time Guidance with Distribution-Shift-Aware Step Sizes}
\author{Philip Nielsen}
\date{Spring 2026}

\begin{document}

% ═══════════════════════════════════════════════════════════════════════════════
\begin{frame}
\titlepage
\end{frame}

% ═══════════════════════════════════════════════════════════════════════════════
\section{Algorithm}

\begin{frame}{Algorithm: MD Diffusion Policy}
\begin{fitcenter}
\begin{algorithmic}[1]
\REQUIRE Policy $\pi_\theta$ (diffusion), critics $Q_{\phi_k}$, replay buffer $\mathcal{D}$
\REQUIRE Energy scale $\alpha$, guidance $\beta$, discount $\gamma$, target update rate $\tau_Q$
\STATE Initialize target network parameters: $\bar{\phi}_k \leftarrow \phi_k$
\FOR{each environment step}
  \STATE \textcolor{gray}{\textit{// Sample action}}
  \STATE $a \leftarrow \textsc{SampleFromTiltedPolicy}(s,\, \pi_\theta,\, \tfrac{1}{2}(Q_{\phi_1}{+}Q_{\phi_2}),\, \alpha,\,\beta)$ \COMMENT{Alg.~2; mean Q for env}
  \STATE \textcolor{gray}{\textit{// Collect experience}}
  \STATE Execute $a$; observe reward $r$, next state $s'$, termination flag $d$
  \STATE $\mathcal{D} \leftarrow \mathcal{D} \cup \{(s, a, r, s', d)\}$
  \STATE \textcolor{gray}{\textit{// Critic update}}
  \STATE Sample minibatch $\{(s_j, a_j, r_j, s'_j, d_j)\}_{j=1}^B \sim \mathcal{D}$
  \STATE $a'_j \leftarrow \textsc{SampleFromTiltedPolicy}(s'_j,\, \pi_\theta,\, \min(Q_{\phi_1},Q_{\phi_2}),\, \alpha,\,\beta)$ \COMMENT{min Q for TD}
  \STATE $y_j \leftarrow r_j + \gamma\,(1{-}d_j)\,\min_k(Q_{\bar{\phi}_k}(s'_j, a'_j))$
  \STATE Update each $Q_{\phi_k}$ by minimizing $\frac{1}{B}\sum_j \textsc{Huber}_{30}(Q_{\phi_k}(s_j, a_j) - y_j)$
  \STATE $\bar{\phi}_k \leftarrow \tau_Q\,\phi_k + (1{-}\tau_Q)\,\bar{\phi}_k$ for $k{=}1,2$ \COMMENT{Polyak averaging}
  \STATE \textcolor{gray}{\textit{// Policy update}}
  \STATE Fit $\pi_\theta(s'_j)$ to $a'_j$ by minimizing diffusion loss
\ENDFOR
\end{algorithmic}
\end{fitcenter}
\end{frame}

% ═══════════════════════════════════════════════════════════════════════════════
\begin{frame}{Q Aggregation Ablation}
\begin{center}
\includegraphics[width=\textwidth,height=0.82\textheight,keepaspectratio]{sweep80/sweep80_ablation_curves_q_agg_sample.png}
\end{center}
\end{frame}

% ═══════════════════════════════════════════════════════════════════════════════
\begin{frame}{Q TD Huber Width Ablation}
\begin{center}
\includegraphics[width=\textwidth,height=0.82\textheight,keepaspectratio]{sweep80/sweep80_ablation_curves_q_td_huber_width.png}
\end{center}
\end{frame}

% ═══════════════════════════════════════════════════════════════════════════════
\begin{frame}{Guidance Strength Ablation}
\begin{center}
\includegraphics[width=\textwidth,height=0.82\textheight,keepaspectratio]{sweep80/sweep80_ablation_curves_tfg_eta.png}
\end{center}
\end{frame}

% ═══════════════════════════════════════════════════════════════════════════════
\begin{frame}{Q Aggregation Comparison Across Matched Configurations}
\small
\begin{center}
\includegraphics[width=\textwidth,height=0.72\textheight,keepaspectratio]{sweep80/sweep80_config_matrix_exp_log_score.png}
\end{center}
Aggregating $Q$ using mean during sampling does not just outperform ``min'' on average. It outperforms in almost every pair of configurations that differ only in the way they aggregate $Q$.
\end{frame}

% ═══════════════════════════════════════════════════════════════════════════════
\begin{frame}{Best Config vs MGMD}
\begin{center}
\includegraphics[width=\textwidth,height=0.82\textheight,keepaspectratio]{sweep80/sweep80_best_vs_mgmd_exp_log_score.png}
\end{center}
\end{frame}

% ═══════════════════════════════════════════════════════════════════════════════
\begin{frame}{Best Config vs Baselines}
\begin{center}
\includegraphics[width=\textwidth,height=0.82\textheight,keepaspectratio]{sweep80/sweep80_best_vs_baselines_exp_log_score.png}
\end{center}
\end{frame}

% ═══════════════════════════════════════════════════════════════════════════════
\begin{frame}{Diffusion Models: Forward and Reverse Processes}
\small
Let $x_0$ denote a clean action drawn from a complex policy distribution. A diffusion model defines a forward noising process
\[
q(x_t\mid x_0)=\mathcal{N}\!\bigl(\sqrt{\bar\alpha_t}\,x_0,\,(1-\bar\alpha_t)I\bigr),\qquad x_t=\sqrt{\bar\alpha_t}\,x_0+\sqrt{1-\bar\alpha_t}\,\epsilon,\quad \epsilon\sim\mathcal{N}(0,I).
\]
As $t$ increases, this smoothly transports a complicated data distribution toward a standard normal distribution.

\medskip
The reverse process starts from $x_T\sim\mathcal{N}(0,I)$ and repeatedly denoises. A network predicts the noise component,
\[
\hat\epsilon_\theta(x_t,s,t)\approx \epsilon.
\]
which gives the Tweedie reconstruction
\[
\hat x_0(x_t,t)=\frac{x_t-\sqrt{1-\bar\alpha_t}\,\hat\epsilon_\theta(x_t,s,t)}{\sqrt{\bar\alpha_t}}.
\]
\end{frame}

% ═══════════════════════════════════════════════════════════════════════════════
\begin{frame}{Energy Guidance and the Noise-Conditional Critic Problem}
\small
Our clean mirror-descent target is the tilted policy
\[
\pi_{\mathrm{new}}(a\mid s)\propto \pi_{\mathrm{old}}(a\mid s)^\alpha\,\exp\!\bigl(\beta\,Q(s,a)\bigr).
\]
At intermediate noise levels, a natural idea is to define a noise-conditional guided predictor,
\[
\hat\epsilon_t^{\mathrm{guided}}(x_t) = \alpha\,\hat\epsilon_\theta(x_t,s,t) - \beta\sqrt{1-\bar\alpha_t}\,\nabla_{x_t}Q_t(s,x_t),
\qquad
Q_t(s,x_t) \approx \mathbb{E}[Q(s,x_0)\mid x_t,s].
\]

\medskip
If $Q_t$ is perfectly trained, then replacing $\hat\epsilon_\theta$ with $\hat\epsilon_t^{\mathrm{guided}}$ in standard diffusion samplers yields
\[
p_t^{\mathrm{tilt}}(x_t\mid s)\propto p_t(x_t\mid s)^\alpha\,\exp\!\bigl(\beta\,Q_t(s,x_t)\bigr),
\]

\medskip
The difficulty is that $Q_t(s,x_t)$ is itself a hard object: it must summarize how clean high-value actions are distributed underneath a noisy sample $x_t$. So although the construction is theoretically exact, the required noisy critic is much harder to learn than an ordinary clean $Q(s,a)$.
\end{frame}

% ═══════════════════════════════════════════════════════════════════════════════
\begin{frame}{Tweedie Guidance: Stronger Signal, No Target Distribution}
\small
Instead of learning a noise-conditional critic, we can guide with a clean critic evaluated at the Tweedie estimate:
\[
\hat\epsilon_t^{\mathrm{Tw}}(x_t) = \alpha\,\hat\epsilon_\theta(x_t,s,t) - \beta\sqrt{1-\bar\alpha_t}\,\nabla_{x_t}Q\bigl(s,\hat x_0(x_t,t)\bigr).
\]
This is attractive in practice because the critic only needs to model clean actions and the guidance signal is typically much stronger, but this changes the level-wise score fields. In general,
\[
\hat\epsilon_t^{\mathrm{Tw}}(x_t) \neq \hat\epsilon_t^{\mathrm{guided}}(x_t),
\]
so Tweedie guidance is no longer the exact reverse process of a prescribed tilted diffusion model. That is, matching the distribution implied by a score field at a particular noise level doesn't imply matching the distribution at the next cleaner noise level under standard diffusion samplers. Still, the sequence has the right endpoints:
\begin{itemize}
  \item at the highest noise level it starts near the Gaussian base distribution,
  \item as $t\to 0$, $\hat x_0(x_t,t)\to x_0$, so the clean composite target $\pi_{\mathrm{new}}(a\mid s)\propto \pi_{\mathrm{old}}(a\mid s)^\alpha e^{\beta Q(s,a)}$ is recovered.
\end{itemize}
This motivates MGMD: use clean-$Q$ Tweedie guidance for a strong practical signal, but correct the induced level-wise distributions so the guided actions are still theoretically sampled from a principled policy update distribution in functional space (Mirror Descent).
\end{frame}

% ═══════════════════════════════════════════════════════════════════════════════
\begin{frame}{Langevin Dynamics on an Energy Landscape}
\small
Suppose we want to sample from a target density of the form
\[
\rho(x) \propto \exp\!\bigl(-E(x)\bigr).
\]
The continuous-time Langevin diffusion is
\[
dX_\tau = -\nabla E(X_\tau)\,d\tau + \sqrt{2}\,dW_\tau,
\]
whose stationary distribution is $\rho(x)$ under standard regularity conditions.

\medskip
In practice we use the Euler--Maruyama discretization,
\[
x^{k+1} = x^k - s\,\nabla E(x^k) + \sqrt{2s}\,\xi^k,
\qquad \xi^k\sim\mathcal{N}(0,I).
\]
This only approaches the desired stationary law in the limit of infinitely many steps and vanishing step size $s\to 0$.

\medskip
So there is a tradeoff: smaller $s$ improves asymptotic accuracy, but slows convergence to the stationary distribution.
\end{frame}

% ═══════════════════════════════════════════════════════════════════════════════
\begin{frame}{Metropolis--Hastings}
\small
Metropolis--Hastings is useful when we already have a Markov chain $q(y\mid x)$ whose stationary distribution is close to the desired sampling distribution.

\medskip
If moves from $x$ to $y$ happen too often relative to the reverse moves, we restrict some of them using an acceptance probability $A(x,y)$ so that detailed balance holds:
\[
\rho(x)\,q(y\mid x)\,A(x,y)
=
\rho(y)\,q(x\mid y)\,A(y,x).
\]
One choice that handles both directions at once is
\[
A(x,y)=\min\!\left\{1,\frac{\rho(y)\,q(x\mid y)}{\rho(x)\,q(y\mid x)}\right\}.
\]

\medskip
Intuition:
\begin{itemize}
  \item if $x\to y$ already happens too frequently, the ratio is $<1$ and we reject some forward moves,
  \item if $x\to y$ happens too rarely, the ratio is $\ge 1$ and the move is accepted with probability $1$.
\end{itemize}
The only extra ingredient MH needs is the ratio of target densities at two points.
\end{frame}

% ═══════════════════════════════════════════════════════════════════════════════
\begin{frame}{Why Energy Parameterization Matters for MALA}
\small
Euler--Maruyama only needs a score field $\nabla \log \rho(x)$, but MH needs the ratio $\rho(y)/\rho(x)$.

 \medskip
 If we write the target in unnormalized energy form,
 \[
 \rho(x)\propto e^{-E(x)},
 \qquad
 \frac{\rho(y)}{\rho(x)} = \exp\bigl(-(E(y)-E(x))\bigr),
 \]
 then the unknown normalizing constant cancels.

 \medskip
This is why MGMD uses an energy-based policy network. Now both ingredients are available:
  \begin{itemize}
    \item gradients for Langevin proposals,
    \item energy differences for MH corrections.
  \end{itemize}
  However, diffusion models are typically trained with MSE loss on $\hat\epsilon_\theta$ so we must determine what $E_\theta(x_t)$ should be given $\hat\epsilon_\theta$.
 \end{frame}

 
 % ═══════════════════════════════════════════════════════════════════════════════
 \begin{frame}{Marginal Noisy Density from the Forward Process}
 \small
 Start from the forward corruption model
 \[
 q(x_t\mid x_0)=\mathcal{N}\!\bigl(\sqrt{\bar\alpha_t}\,x_0,(1-\bar\alpha_t)I\bigr).
 \]
 If clean data are distributed as $x_0\sim p_{\mathrm{data}}(x_0)$, then the noisy marginal is obtained by integrating out $x_0$:
 \[
 p_t(x_t)=\int p_{\mathrm{data}}(x_0)\,q(x_t\mid x_0)\,dx_0.
 \]

 \medskip
 Since any positive density can be written in unnormalized energy form, define
 \[
 p_t(x_t)=A_t\exp\!\bigl(-E_t(x_t)\bigr),
 \]
 where $A_t$ is a normalization constant that depends on $t$ but not on $x_t$.

 \medskip
 Therefore,
 \[
 \nabla_{x_t}\log p_t(x_t)=-\nabla_{x_t}E_t(x_t).
 \]
 So if we can express the score $\nabla_{x_t}\log p_t(x_t)$, we immediately get the energy gradient used by Langevin and MALA.
 \end{frame}

 % ═══════════════════════════════════════════════════════════════════════════════
 \begin{frame}{Differentiate the Marginal via the Posterior over $x_0$}
 \small
 Differentiate the marginal under the integral sign:
 \[
 \nabla_{x_t}p_t(x_t)
 = \int p_{\mathrm{data}}(x_0)\,\nabla_{x_t}q(x_t\mid x_0)\,dx_0.
 \]
 Use $\nabla q = q\,\nabla \log q$ inside the integral:
 \[
 \nabla_{x_t}p_t(x_t)
 = \int p_{\mathrm{data}}(x_0)\,q(x_t\mid x_0)\,\nabla_{x_t}\log q(x_t\mid x_0)\,dx_0.
 \]
 Divide by $p_t(x_t)$ and apply Bayes' rule:
 \[
 \begin{aligned}
 \nabla_{x_t}\log p_t(x_t)
 &= \frac{1}{p_t(x_t)}\int p_{\mathrm{data}}(x_0)\,q(x_t\mid x_0)\,\nabla_{x_t}\log q(x_t\mid x_0)\,dx_0 \\
 &= \int p(x_0\mid x_t)\,\nabla_{x_t}\log q(x_t\mid x_0)\,dx_0 \\
 &= \mathbb{E}_{x_0\mid x_t}\!\left[\nabla_{x_t}\log q(x_t\mid x_0)\right].
 \end{aligned}
 \]
 So the noisy score is the posterior average of the Gaussian score coming from the forward kernel.
 \end{frame}

 % ═══════════════════════════════════════════════════════════════════════════════
 \begin{frame}{From the Gaussian Score to Tweedie's Formula}
 \small
 For the Gaussian forward kernel,
 \[
 \log q(x_t\mid x_0)=\mathrm{const}-\frac{\|x_t-\sqrt{\bar\alpha_t}x_0\|^2}{2(1-\bar\alpha_t)},
 \]
 so
 \[
 \nabla_{x_t}\log q(x_t\mid x_0)
 = -\frac{x_t-\sqrt{\bar\alpha_t}x_0}{1-\bar\alpha_t}.
 \]
 Taking the posterior expectation gives
 \[
 \nabla_{x_t}\log p_t(x_t)
 = -\frac{x_t-\sqrt{\bar\alpha_t}\,\mathbb{E}[x_0\mid x_t]}{1-\bar\alpha_t}.
 \]
 Therefore,
 \[
 \nabla_{x_t}E_t(x_t)
 = \frac{x_t-\sqrt{\bar\alpha_t}\,\mathbb{E}[x_0\mid x_t]}{1-\bar\alpha_t},
 \qquad
 \mathbb{E}[x_0\mid x_t]
 = \frac{x_t-(1-\bar\alpha_t)\nabla_{x_t}E_t(x_t)}{\sqrt{\bar\alpha_t}}.
 \]
 Finally, since $x_t=\sqrt{\bar\alpha_t}x_0+\sqrt{1-\bar\alpha_t}\,\epsilon$,
 \[
 \mathbb{E}[\epsilon\mid x_t]
 = \frac{x_t-\sqrt{\bar\alpha_t}\,\mathbb{E}[x_0\mid x_t]}{\sqrt{1-\bar\alpha_t}}
 = \sqrt{1-\bar\alpha_t}\,\nabla_{x_t}E_t(x_t).
 \]
 This is the link between the energy gradient, Tweedie's conditional mean, and the diffusion model's noise prediction.
 \end{frame}

 % ═══════════════════════════════════════════════════════════════════════════════
 \begin{frame}{How This Implies We Should Parameterize the Model}
 \small
 The identity from the previous slide suggests parameterizing the denoising noise field as
 \[
 \hat\epsilon_\theta(s,x_t,t)\coloneqq\sqrt{1-\bar\alpha_t}\,\nabla_{x_t}E_\theta(s,x_t,t).
 \]

 In practice, $\hat\epsilon_\theta$ is still trained with the standard diffusion noise-prediction loss
 \[
 \mathcal{L}_{\mathrm{diff}}(\theta)=\mathbb{E}_{x_0,\epsilon,t}\left[\left\lVert \hat\epsilon_\theta\!\left(s,\sqrt{\bar\alpha_t}x_0+\sqrt{1-\bar\alpha_t}\,\epsilon,t\right)-\epsilon\right\rVert^2\right].
 \]

 define the composite guided energy
 \[
 E_t^{\mathrm{guided}}(x_t)=\alpha\,E_\theta(s,x_t,t)-\beta\,Q\bigl(s,\hat x_0(x_t,t)\bigr).
 \]
 Therefore,
 \[
 \sqrt{1-\bar\alpha_t}\,\nabla_{x_t}E_t^{\mathrm{guided}}(x_t)
 = \alpha\,\hat\epsilon_\theta(s,x_t,t)-\beta\sqrt{1-\bar\alpha_t}\,\nabla_{x_t}Q\bigl(s,\hat x_0(x_t,t)\bigr)
 = \hat\epsilon_t^{\mathrm{Tw}}(x_t).
 \]
 \end{frame}

 % ═══════════════════════════════════════════════════════════════════════════════
 \begin{frame}{From Langevin $+$ MH to MALA at Each Diffusion Level}
 \small
 MGMD then targets the level-wise density
 \[
 p_t(x\mid s)\propto \exp\!\bigl(-E_t^{\mathrm{guided}}(x)\bigr).
 \]
 Using Langevin, the proposal kernel is
 \[
 q_t(x'\mid x)=\mathcal{N}\!\bigl(\mu(x),\; 2s_t I\bigr),\qquad \mu(x)=x-s_t\nabla E_t^{\mathrm{guided}}(x),
 \]
 where $s_t$ is the MALA step size at that noise level. Applying MH to this proposal gives MALA, with acceptance probability
 \[
 \begin{aligned}
A(x,x')&=\min\!\left\{1,\exp\!\left(E_t^{\mathrm{guided}}(x)-E_t^{\mathrm{guided}}(x')\right.
 + \left.\frac{\|x'-\mu(x)\|^2-\|x-\mu(x')\|^2}{4s_t}\right)\right\}.
\end{aligned}
\]
So MALA removes the need to send $s_t\to 0$ for correctness. We can take reasonably large guided steps while still converging to the desired level-wise target distribution in the limit of infinite steps.
\end{frame}

% ═══════════════════════════════════════════════════════════════════════════════
\begin{frame}{Why Parameterize by $(\alpha,T)$ Instead of $(\alpha,\beta)$?}
\small
Start from the composite update
\[
\pi_{k+1}(a\mid s)\propto \pi_k(a\mid s)^\alpha\,\exp\!\bigl(\beta\,Q(s,a)\bigr).
\]
Set
\[
\beta=\frac{1-\alpha}{T}.
\]
So in log-density form,
\[
\log \pi_{k+1}(a\mid s)
= \alpha\log \pi_k(a\mid s) + (1-\alpha)\frac{Q(s,a)}{T} - \log Z_k(s).
\]
At a fixed point,
\[
\pi_*(a\mid s)\propto \exp\!\bigl(Q(s,a)/T\bigr).
\]
\begin{itemize}
  \item $\alpha$: memory / relaxation rate,
  \item $T$: asymptotic Boltzmann temperature (analogous to SAC's entropy temperature $\alpha$, since SAC has $\pi^*(a\mid s)\propto e^{Q(s,a)/\alpha}$),
  \item sweeping $(\alpha,T)$ separates update timescale from exploitation sharpness.
\end{itemize}
\end{frame}

% ═══════════════════════════════════════════════════════════════════════════════
\begin{frame}{Algorithm 2: \textsc{SampleFromTiltedPolicy} (MALA $+$ DDIM)}
\small
\begin{fitcenter}
\begin{algorithmic}[1]
\REQUIRE State $s$, energy policy $E_\theta$, aggregated critic $\bar Q(s,a)$, energy scale $\alpha$, guidance $\beta$
\REQUIRE Schedule $\{\bar\alpha_t\}$, per-level step sizes $\{s_t\}$.
\REQUIRE Q-eval clip $r_Q$, MALA steps $M$
\STATE $\hat\epsilon_\theta(x,t)\coloneqq\sqrt{1{-}\bar\alpha_t}\,\nabla_x E_\theta(s,x,t)$
\STATE $\hat x_{0,r_Q}(x,t)\coloneqq \mathrm{clip}\!\left(\tfrac{x-\sqrt{1{-}\bar\alpha_t}\,\hat\epsilon_\theta(x,t)}{\sqrt{\bar\alpha_t}},\;{-}r_Q,\;r_Q\right)$
\STATE $E_t^{\mathrm{guided}}(x)\coloneqq \alpha\,E_\theta(s,x,t)-\beta\,\bar Q\bigl(s,\hat x_{0,r_Q}(x,t)\bigr)$
\STATE Sample $x_T\sim\mathcal{N}(0,I)$
\FOR{$t=T,\dots,1$}
  \STATE \textcolor{gray}{\textit{// corrector step}}
  \FOR{$m=1,\dots,M$}
    \STATE $\mu\leftarrow x_t-s_t\nabla E_t^{\mathrm{guided}}(x_t)$;\quad Sample proposal $x'\sim\mathcal{N}(\mu,2s_tI)$ \COMMENT{Langevin proposal}
    \STATE $\mu'\leftarrow x'-s_t\nabla E_t^{\mathrm{guided}}(x')$ \COMMENT{reverse proposal mean}
    \STATE $A\leftarrow\min\!\left\{1,\exp\!\bigl(E_t^{\mathrm{guided}}(x_t)-E_t^{\mathrm{guided}}(x')\bigr)\tfrac{\mathcal{N}(x_t;\mu',2s_tI)}{\mathcal{N}(x';\mu,2s_tI)}\right\}$ \COMMENT{MH accept prob.}
    \STATE Set $x_t\leftarrow x'$  with probability $A$ \COMMENT{accept/reject}
    \STATE Update $s_t$ online using the current accept/reject outcome \COMMENT{adapt toward target rate}
  \ENDFOR
  \STATE \textcolor{gray}{\textit{// predictor step}}
  \STATE $\hat\epsilon_g\leftarrow\sqrt{1{-}\bar\alpha_t}\,\nabla_{x_t}E_t^{\mathrm{guided}}(x_t)$
  \STATE $x_{t-1}\leftarrow\sqrt{\bar\alpha_{t-1}/\bar\alpha_t}\,x_t + \left(\sqrt{1{-}\bar\alpha_{t-1}}-\sqrt{\bar\alpha_{t-1}(1{-}\bar\alpha_t)/\bar\alpha_t}\right)\hat\epsilon_g$ \COMMENT{standard DDIM denoising step}
\ENDFOR
\STATE \textbf{return} $\mathrm{clip}(x_0,{-}1,1)$, updated $\{s_t\}$
\end{algorithmic}
\end{fitcenter}
\end{frame}


\end{document}
